\chapter{話し方}

この講義は、話すことを、構造と技法と帰結をもつひとつの技芸として扱う。始まりは小手先の助言ではない。まず切迫感が提示され、その切迫感が伝達の質をめぐる一つの公式へと絞り込まれ、それからようやく、話をどう始めるか、どうすれば聞き手を離れさせずにすむか、黒板や小道具をどう使うか、なぜスライドはしばしば失敗するのか、ジョブトークでは何をほとんど即座に示さねばならないのか、そして話し手はいかに締めくくるべきかという仕掛けが、順を追って解きほぐされてゆく。この順序は教育的である。まず、ここには学ぶべきことがあるのだと納得させられ、そののち、その知識が実際にはどのような姿をとるのかを見せられるのである。

\section{冒頭のフック、コミュニケーションの公式、そして約束}

Winston は、強烈な類比から語り起こす。彼によれば、Uniform Code of Military Justice は、兵士に武器を持たせず戦場へ送り出した士官を処罰する。学生に対しても、それに対応する保護があってしかるべきだ。なぜなら、学生もまた、伝える能力を持たぬまま人生へ送り出されるべきではないからである。問題は文学的な洗練ではない。問題は、アイデアはひとりでに世の中を渡ってはくれない、ということにある。

ついで彼は、その優先順位を、望みうるかぎり率直に述べる。

\begin{equation}
\text{話すこと} \succ \text{書くこと} \succ \text{アイデアの質}.
\end{equation}

ここで \(\succ\) は、彼の言う ``その順で'' を編集上手短に記したものにすぎない。この主張は、アイデアが本来的に伝え方より重要でない、という意味ではない。むしろ、世の中での成否は、こうした事柄がどの順番で失敗し、あるいは成功するかによって、しばしばふるいにかけられるのだ、という意味である。

そこから講義は、最初の形式的対象へと絞り込まれる。後年の黒板に残った断片ではなく、確実な transcript の行に基づけば、この導入の公式は、慎重を期して標準的な形で次のように書ける。

\begin{equation}
Q = f(K,P,T),
\end{equation}

ここで \(Q\) は話すこと、あるいは書くことの質を表し、\(K\) は知識、\(P\) は練習、そして \(T\) は生得的な才能である。

\begin{remark}
のちの教室の映像では、自転車の車輪のある場面にも、この公式の背景の残滓、すなわち \(\text{Quality}=f(K,P,\ldots)\) のような部分形しか見えていない。したがって、ここでの三変数の完全な形は、スクリーンショットだけを完全に書き写したものではなく、transcript に基づいて復元したものである。
\end{remark}

講義はすぐに、この公式の読み方を教える。それは引数に関して対称ではない。Winston は、\(T\) は ``非常に小さい'' と言う。彼自身の黒板記法ではなく、その意味を忠実に要約する簡潔な数学的略記で言えば、こう書ける。

\begin{equation}
T \ll K,P.
\end{equation}

ここでの要点は、単に ``練習'' ではなく、``知識'' という語である。Winston はきわめて明確に、ほんとうに大事なのは、あなたが何を知っているかだと言う。

そして彼は、この公式を抽象的な議論ではなく、一つの話によって試す。Mary Lou Retton は並外れたアスリートだが、スキーの初心者である。Winston は運動能力において彼女よりはるかに恵まれていないが、それでもよりよいスキーヤーである。なぜか。その課題に関しては、彼のほうに必要な \(K\) と \(P\) があり、彼女には主として \(T\) があるだけだからである。この逸話は、公式のまわりに置かれた飾りではない。この公式を本気で受け取るべきだという最初の証明なのである。

\subsection{質問 \& 答え}

\paragraph{質問。} なぜ \(T\) は小さいのか。

\paragraph{答え。} 才能は、それだけでは、まだ課題に即した能力ではないからである。スキーの例は、その点をこれ以上ないほど鋭く示している。Retton には莫大な一般的運動才能があるが、Winston にはスキーについての知識とスキーの練習がある。だから、スキー板の上では彼のほうが優れている。講義が求めているのは、この構造をそのまま話すことへ移し替えることである。生来の資質が中程度でも、よい手本と習慣と準備を備えた人は、より自然な才能を持ちながら訓練の足りない話し手を凌ぐことができる。

この逸話から講義は、ただちに「約束」へ移る。もし公式が正しいなら、上達は神秘ではない。聴衆には、話すための技法の武器庫が約束され、Winston はさらに重要な但し書きを添える。その過程は非線形である、と。そうした技法のうちの一つ、おそらくはただ一つが、誰かに職をもたらす決め手になるかもしれない。この一言は投げ捨てられた文句ではない。理論から、このあとの一時間全体へ橋を架ける、動機づけの継ぎ目なのである。

\section{受講の規律と、どう始めるか}

この約束を与えてはじめて、Winston は規律を課す。ノートパソコン禁止、携帯電話禁止という規則は、礼儀作法の問題として提示されるのではない。認知の事実として提示されるのである。人間には、彼の言うところでは、言語処理装置が一つしかない。もしその処理装置がメールを読み、ウェブを見ているなら、それは話を聞いてはいない。しかも害は社会的でもある。周囲の聞き手は気を散らされ、開いたノートパソコンが視界に残っていると、話し手の出来さえ悪くなる。

これはこの講義を貫く中心的な考えの一つである。のちには、文字だらけのスライドに反対する最も深い議論として戻ってくる。ここでそれが果たしている役割は別である。本論が始まる前に、まず伝達の回路を澄ませるのである。

次の一手も同じく意図的だ。``最初に話すべきことは'' と Winston は言う。``どう始めるかだ。'' そしてすぐに、よくある助言を退ける。ジョークで始めるべきだ、という助言である。理由は、ジョークが真面目な話にふさわしくないからではない。理由はタイミングにある。話の初めには、部屋全体がまだ話し手の声、速度、リズムに順応している最中だからだ。聴衆はまだ同調していない。

\subsection{質問 \& 答え}

\paragraph{質問。} なぜジョークで始めてはいけないのか。

\paragraph{答え。} ジョークには落ち着いた聴衆が必要だが、話の冒頭にあるのは落ち着いていない聴衆だからである。人はまだ入室中であり、まだ物をしまっており、まだ話し手の声のパラメータに身体を合わせている。したがってジョークは、聴衆がそれを受け取る準備を整える前にやって来てしまう。むしろ冒頭にふさわしいのは、力を与える約束である。すなわち、この一時間の終わりには、初めには知らなかった何を聴衆が知ることになるのかを、はっきり述べることである。

これはこの講義の、もっとも気の利いた自己反照的な特徴の一つである。Winston は、力を与える約束を勧めているだけではない。すでにそれを自ら実行している。講義は切迫感で始まり、ついで公式が示され、逸話による証明が与えられ、いまや使える技法への約束に至った。話そのものが、すでに自分自身の実例として用いられているのである。

\section{四つのヒューリスティック}

聴衆が、ここには確かな知識があるのだと納得したところで、Winston は、話すときつねに念頭に置いている四つの見本的ヒューリスティックを挙げる。主題を巡回すること、アイデアのまわりに柵を立てること、言葉による句読点を打つこと、そして質問することである。

\paragraph{巡回する。} 最初のヒューリスティックは、その主題のまわりを一度回り、そしてもう一度回ることだ。Winston は昔からの助言を思い出させる。これから何を話すかを言い、実際にそれを話し、そして三度目にもう一度それを言うのである。彼の正当化は、聞き手が鈍いからではない。聞き手は断続的だからである。どんな講義であれ、ある瞬間には部屋のおよそ \(20\%\) が、彼の言うところでは、どうしても霧の中にいる。

このヒューリスティックは、小さな編集上のモデルを誘う。もし \(p_{\mathrm{fog}}\) が、聞き手が一回の通過を聞き逃す確率を表すなら、Winston の \(20\%\) という見積もりは次を与える。

\begin{align}
p_{\mathrm{fog}} &\approx 0.2,\\
\mathrm{Prob}(\text{三回とも聞き逃す}) &\approx p_{\mathrm{fog}}^3 = 0.2^3 = 0.008,\\
\mathrm{Prob}(\text{少なくとも一回は受け取る}) &\approx 1-p_{\mathrm{fog}}^3 = 0.992.
\end{align}

\begin{remark}
これは Winston 自身の計算ではなく、われわれによる説明上の補注である。講義の本当の要点は確率論ではなく実践にある。反復は、話の中へ再入場する機会を複数つくるのである。
\end{remark}

\subsection{質問 \& 答え}

\paragraph{質問。} なぜ三回言うのか。

\paragraph{答え。} 話すことは時間の中で起こり、聴衆は完全に静止してはいてくれないからである。一度きりの通過は脆い。二度目、三度目の通過は、とりわけ少しずつ異なる側面から同じ点に近づくとき、伝達を頑丈なものにする。したがって ``巡回'' は、空疎な重複ではない。注意の漂流に対する工学的な応答なのである。

\paragraph{柵を立てる。} 第二のヒューリスティックは、対比による定義である。あるアイデアが何であるかを聴衆に分かってほしいなら、それが何でないかも見えるようにしてやるべきだ。Winston の簡単な例はアーチである。これはアーチだ、これはアーチではない、あれもアーチではない。技術的な形でも同様に有効である。私のアルゴリズムは Jones のアルゴリズムに似ているが、ただし彼のものは指数的で、私のものは線形だ。概念は、そのまわりに柵を築くとき、輪郭を獲得する。

\paragraph{言葉による句読点。} 第三のヒューリスティックも、同じ注意診断から出てくる。人が途中でバスから落ちるなら、話のなかに、そこからまた乗り直せる継ぎ目が必要になる。明示的な列挙、要約、転換句は、その継ぎ目を提供する。Winston 自身の講義も、第一のアイデアから第二へ進み、いま第三に来たと告げることで、この方法を実演している。

\paragraph{質問する。} 第四のヒューリスティックは、もっとも直接的に聴衆を引き戻す装置である。質問は受動的な受容を中断し、部屋全体をひとつの問題へ向かって身構えさせる。しかしここでも講義は、節度ある規則を付け加える。十分な長さ、待たなければならない。

\begin{equation}
t_{\mathrm{wait}} \approx 7\,\mathrm{s}.
\end{equation}

七秒というのは、話し手には永遠のように感じられる。まさにそのために、この規則が必要なのだ。部屋には考える時間が要る。同時に、その質問は慎重に選ばれねばならない。あまりに自明で、答えるのが気まずくなるほどであってもいけないし、あまりに難しくて、誰も何も言えなくなるほどであってもいけない。

このブロックの終わりで、Winston は示唆的なことを口にする。もし聴衆が、いまや話すことには ``学ぶべき何か'' があると納得しているなら、それだけですでに彼は成功しているのだ、と。ここが講義の折れ目である。はじめてここで、彼は視野を広げる。

\section{議論の一部としての時間と場所}

Winston によれば、次の議題は時間と場所である。この転換は明示的で、講義はこの拡大が当然の報いとして得られたものだと感じさせようとしている。まず私たちは、話の内側の力学を扱った。つぎに、その力学が働きうる条件を考えるのである。

講義に好ましい時間は午前 \(11\) 時である。Winston の理由は実際的で、身体に根ざしている。人が目覚めており、食後すぐではなく、まだ疲労に落ち込んでもいない時間だからだ。望ましい部屋は明るい。薄暗い部屋は眠りを合図するからである。また、あらかじめ下見しておくべきでもある。つまり、その場での不意打ちを最小限にするため、事前に点検しておくべきだということだ。さらに、部屋にはそれなりに人が入っているべきである。半分空いたホールは、それだけで気勢をそぐ社会的な芝居を生んでしまう。

こうした点は、講義中もっとも優れた黒板フレームの一つに、視覚的なかたちで残っている。このスクリーンショットが貴重なのは、そこに導出が書かれているからではなく、黒板の建築を保存しているからである。

\begin{figure}[h]
\centering
\includegraphics[width=0.95\textwidth]{lecture_11_figure_03.png}
\end{figure}

左側には \emph{例} という見出しがあり、四つのヒューリスティックが並んでいる。``巡回'' のそばの輪は形式的な図ではないが、反復通過という考えを視覚的に符号化する有益なしるしになっている。右側には \emph{時間 \& 場所} という見出しがあり、長いうねる線で区切られている。この黒板はすでに、Winston がよい黒板の条件として述べることを実行している。すなわち、話を、単に言葉によってではなく、空間的にも組織しているのである。

整えた再構成図にすると、その組織はさらに読み取りやすくなる。右側の正確な文言は、そのフレームだけからでは完全には確定しないので、以下の項目は transcript によって補いながら、黒板全体の幾何学は保ってある。

\begin{figure}[h]
\centering
\begin{tikzpicture}[x=1cm,y=1cm]
\draw[thick] (0,0) rectangle (13.5,6.8);

\node[anchor=west] at (0.6,6.0) {\Large 例};
\node[anchor=west] at (0.8,5.0) {巡回};
\draw[thick] (3.0,5.0) ellipse (0.45 and 0.95);
\draw[thick] (3.45,4.95) ellipse (0.45 and 0.95);
\draw[->,thick] (4.0,5.0) -- (4.5,5.0);

\node[anchor=west] at (0.8,4.0) {柵を立てる};
\node[anchor=west] at (0.8,3.0) {言葉による句読点};
\node[anchor=west] at (0.8,2.0) {質問};

\node[anchor=west] at (6.0,6.0) {\Large 時間 \& 場所};
\draw[thick] plot[smooth] coordinates {(5.8,4.85) (6.8,4.95) (7.8,4.8) (8.8,4.95) (9.8,4.75) (10.8,4.9) (11.8,4.7) (12.8,4.8)};
\node[anchor=west] at (6.5,4.1) {午前11時};
\node[anchor=west] at (6.5,3.3) {明るい};
\node[anchor=west] at (6.5,2.5) {下見済み};
\node[anchor=west] at (6.5,1.7) {人が入っている};
\end{tikzpicture}
\end{figure}

この講義がほんとうに言っているのは、時間と場所は単なる段取り上の細目ではない、ということだ。それらは注意の力学を左右する変数なのである。

\section{黒板と小道具: なぜ物理的な媒体が効くのか}

時間と場所を論じたのち、Winston はようやく道具へ移る。この順序が重要である。道具は真空のなかで選ばれるのではない。部屋のなかで、時と目的に応じて選ばれるのである。

基本的な区別は単純だ。黒板が適切な道具なのは、知らせること、教えること、講義することが目的である場合だ。スライドは、会議発表やジョブトークのように、アイデアを素早く露出させるのに向いている。そこから Winston は、本当の教育において黒板が好ましい理由を三つ挙げる。

第一に、図像的な質がある。黒板では、まとめ、分け、下線を引き、スケッチできる。第二に、速度に関する性質がある。書く速さは、おおむね聴衆が吸収できる速さに近い。スライドは思考を追い越してしまうことがある。第三に、黒板は話し手の手に向かう先を与えてくれる。これは冗談以上のことだ。人前で話すぎこちなさの多くは、身体への突然の自意識から生まれるが、黒板はその身体に、正当な仕事を与えてくれる。

ここから講義は、小道具へと視野を広げる。導きの原理は、ここでもまた理論に先立つ実例である。Ibsen の \emph{Hedda Gabler} におけるストーブが効くのは、それがそこに置かれ、緊張を集め、しだいに避けがたいものになってゆくからである。小道具が記憶に残るのは、単に説明されるのではなく、舞台にのせられているからだ。

同じ論理は教室にも現れる。Papert の自転車の車輪は飾りではない。ひと筋の推論を運ぶための装置である。

\begin{figure}[h]
\centering
\includegraphics[width=0.78\textwidth]{lecture_11_figure_04.png}
\end{figure}

この図が重要なのは、主として物理的な実演を示しているからである。車輪はまっすぐ立ち、ハブは講師の手にあり、小道具全体がはっきり見えている。背景には、導入の黒板公式の部分的な残り、たとえば \(\text{Quality}=f(K,P,\ldots)\) のようなものや、先ほどのストーブのスケッチもまだ見て取れる。しかし、この図の数学的な重みを担っているのは、そうした残骸ではなく、小道具そのものである。

\subsection{質問 \& 答え}

\paragraph{質問。} 回転している車輪は、どちらへ進むのか。

\paragraph{答え。} Winston の要点は、この実演を考えるのに、暗記した手の規則はよい方法ではないということである。彼は修辞的に、多くの機械系の技術者でさえ成功率はせいぜい \(50\%\) ほどだ、と言う。つまり、その規則はコイン投げと同じものになってしまっている。よりよい方法は、問題の表現を変えることである。

\begin{enumerate}
\item 車輪を回し、小さなリムの一片にダクトテープで印をつける。
\item ひとまず、車輪全体として考えることを拒む。
\item 印をつけた部分が上へ回ってくるのを見る。
\item その部分に向けて、息のひと吹き、あるいは局所的に等価な押しを加える。
\item その小さな一片が、いま何をしたがるかを問う。
\item ついで、次の一片が同じ位置に来るのを見る。
\item すると、同じ局所的推論がもう一度適用される。
\item したがって、車輪全体は対応する方向へ回り続けねばならない。
\end{enumerate}

ここでの達成は、完全なジャイロスコープ解析ではない。講義はそのような解析を黒板に書かないし、あたかもそこにあったかのように持ち込むべきでもない。ここでの達成は教育的なものである。この小道具は、暗記された記号操作を、より信頼できる像へと置き換える局所部分の議論を可能にする。

第二の小道具の例は、Alan Lazarus による、エネルギー保存を示す振り子球の実演である。重い球を鼻の高さまで持ち上げ、静止状態から放すと、それは向こうへ振れてゆき、ふたたび戻ってきて、実演者の鼻に触れはするが、打ち砕きはしない。関係する方程式は、よく知られた次のものである。

\begin{equation}
E = K_{\mathrm{kin}} + V_{\mathrm{pot}} = \text{一定}.
\end{equation}

講義はこれを黒板で導出しないし、その必要もない。要点は、この実演が、その法則を身体的に記憶に残るものにする、ということである。本当の危険は、Winston の言うところでは、保存則が破れることではなく、ただ手を離すのでなく人間がつい押してしまう、その傾向にある。

ストーブ、車輪、振り子を見せたあとではじめて、Winston は一般的な説明を与える。彼の言う ``lunatic fringe'' 理論によれば、チョークと小道具が効くのは、共感的なミラーリングが働くからである。手が書くのを見たり、球が振れるのを見たりすると、私たちはほとんど自分で書いたり振れたりしているかのように感じる。スライドや静止画は、同じ身体的参与を生まない。ここでもまた講義は、自らの方法に忠実である。実例が先、説明はあとである。

\section{スライド、特殊な場合、そしてジョブトークの閾値}

Winston によれば、まだ論じるべき道具がもう一つある。スライドである。先ほどの区別は、いまや黒板と小道具が確立されたのち、さらに強く反復される。スライドは、教えるためではなく、アイデアを露出させるためのものだ。

攻撃は、ほとんど例外を認めない法則として述べられる。人はいつもスライドを多くしすぎ、そのスライドにはいつも言葉が多すぎる。そこから講義は、スライドの罪状一覧となる。スライドを声に出して読んではならない。話し手がスライドから離れすぎて、聴衆に画面と話し手のあいだで視線のテニスを強いてはならない。背景の余計な飾り、無意味なロゴ、話し手がすでに声に出して言っている題名は削ぎ落とすべきである。

より深い議論は、冒頭の「一つしかない言語処理装置」という主張と同じである。文字の密なスライドは、聴衆を読書へと引き込み、部屋全体が読み始めた瞬間、それはもう聞いてはいない。Winston は実験によってこの点を補強する。内容の一部が話され、一部がスライドに現れるとき、学生はスライド上の内容のほうをよく覚える。あとで学生の一人が、話し手が話していたのが気を散らした、と言う。この冗談が効くのは、その原理が真実だからである。

そして講義は助言をさらに鋭くする。スライドには普通、ごく少数の語しか載せるべきではない。もし重要な歴史的テクストがあるなら、話し手は明示的に、聴衆に読む時間を与えるべきである。そして興味深い例外が一つある。Winston の \emph{hapax legomenon} スライド、すなわち、複雑さそのものを見せることだけが眼目であるときにかぎって、ただ一度だけ許される、あの途方もなく密なスライドである。だが一度は一度である。

タイポグラフィの規則は実際的だ。

\begin{equation}
40\text{--}50\,\mathrm{pt}\ \text{なら安全}, \qquad 35\,\mathrm{pt}\ \text{ではすでに小さすぎるほうへ流れ始めている。}
\end{equation}

講義の要点は、\(35\) ポイントでは読みにくいかもしれない、ということだけではない。話し手が \(35\) ポイントに手を伸ばすのは、スライドにあまりに多くの言葉を押し込もうとしているときだ、ということなのである。

Winston はさらに別の罪も挙げる。話し手を聴衆から背けてしまうレーザーポインタ。スライドを印刷して机に並べてみれば診断できる、重すぎるデッキ。空気も余白もイメージもないデッキ。積み重なる教訓は、スライドは話された推論の線への薬味であるべきで、壁の上で別の講義を競わせてはならない、ということだ。

道具から講義は、特殊な場合へ移る。第一は informing、すなわち Winston がいままさに行っている種類の話である。ここで力を与える約束が戻ってくるが、新たな強調が加わる。人を鼓舞するには、情報を構造化する以上のことをしなければならない。学生に、君にもできると言ってやってもよいし、問題を新しい見方で見えるようにしてやってもよいし、ただ主題への本物の情熱を示してもよい。資源配分の話がその例である。まず私たちは、壊滅的な総当たり計算を見せられる。ついで、このあと五十分で説明される小さな変更によって、太陽系の寿命に匹敵する計算を数秒に縮められる、と約束される。この約束は説明的であると同時に、胸を躍らせるものでもある。

つぎに、教員がよく自分たちの仕事だと言うことについて、小さな脇道が入る。人に考え方を教えることだ、と。Winston は、そうした主張が当然に呼び寄せる次の問いを発する。どうやってか。彼の答えは、人間は物語を語る動物だ、というものである。したがって、思考を教えるとは、物語を与え、物語についての問いを与え、分析の仕組みを与え、物語を組み合わせる方法を与え、その信頼性を評価する方法を与えることなのである。この脇道が重要なのは、プレゼンテーションを、説得だけでなく教育にも結びつけているからだ。

口頭試問は、さらに二つの助言を与えてくれる。受験者は自分の仕事を位置づけなければならない。つまり、それをより大きな分野の中に置き、その問題がなぜ重要かを説明しなければならない。また、練習は本当の練習でなければならない。すでに内容を知っていて、欠けた説明を黙って補ってくれる人々の前でのリハーサルではだめなのだ。

\subsection{質問 \& 答え}

\paragraph{質問。} ジョブトークでは、どれくらい時間があるのか。

\paragraph{答え。} およそ五分である。

\begin{equation}
t_{\mathrm{job}} \approx 5\,\mathrm{min}.
\end{equation}

その時間幅のなかで、Winston によれば、二つのことを確立しなければならない。聴衆は、その候補者に vision があることを見なければならず、同時に、その候補者が何かをすでに成し遂げていることも見なければならない。vision とは、誰かが気にかけるに足る問題があり、そのアプローチに本当に新しいものがあるということである。``何かを成し遂げている'' とは、単なる志ではなく、具体的な前進があるということだ。

講義はこれを、推奨という形で具体化する。話し手は、その問題を解くために必要な段階を列挙し、そのうちどれが実際に遂行されたかを示し、最後に contributions スライドでその段階を映し返すべきである。こうして、技術的な話の簡潔な構造が浮かび上がる。

\begin{enumerate}
\item 気にかける価値のある問題を述べる。
\item そのアプローチの何が新しいかを明らかにする。
\item 現実の前進に必要となる段階を列挙する。
\item そのうちのいくつかが実際に達成されていることを示す。
\item 最後に、先のリストを映し返す contributions スライドで締めくくる。
\end{enumerate}

これは単なる作法ではない。見知らぬ人がごく短時間のうちに、その話し手が孤立した事実の寄せ集めではなく、研究プログラムを持っているかどうかを、どのように判断するかについての理論なのである。

\section{記憶されることと、終わり方を知ること}

講義の最後の運動は、ジョブトークのあとから始まる。職を得たあとでもなお、Winston は言う。自分の仕事の価値によって認識されることを考えねばならない。そのことを彼は、Julia Child の逸話で舞台化する。彼女のもとには、自分の人生が変わったと言いに来る人がひっきりなしに現れる。有名であるのは楽しいかと問われたときの彼女の答えは、乾いている。人はそれにも慣れるものだ、と。Winston の返答はそれより示唆的だ。人は、無視されることには決して慣れない。この教訓は、パッケージングは虚栄ではない、ということである。それは、アイデアにこの世界で公正な機会を与えるための一部なのだ。

そこから彼は、仕事が記憶に残るようになる五部構造として、Winston's Star を導入する。

\begin{itemize}
\item symbol,
\item slogan,
\item surprise,
\item salient idea,
\item story.
\end{itemize}

あのアーチ学習の古い例が、この星を展開してみせる。アーチそれ自体が symbol である。``One-shot learning'' が slogan である。十分に賢い学習者であれば、一つの例からでも確かなことを学べる、というのが surprise である。salient idea は、もう少しで届かなかったその近傍にある。単なる否定例ではなく、概念が何を本当に要求しているかを明らかにする、ほとんど例であるようなものだ。そして最後に、すべては、そのシステムがどう働き、なぜ重要なのかという story に結びつけられなければならない。Winston は、ここでの ``salient'' は、抽象的に ``重要'' という意味ではないと強調する。それは、記憶に残るほど十分に突き出しているアイデア、という意味なのである。

記憶の話のあとに、ようやく終結が来る。ここで講義は、悪い終わり方を演じてみせる。協力者の巨大な一覧で終わってはならない。そんなことをすれば、聴衆がまさに話し手のしたことを見定めようとしているその瞬間に、話し手自身を小さくしてしまうからだ。最後の画面を空疎な余白に費やしてはならない。ほんとうに重要なのが contribution なら、``conclusions'' で終わってもならない。

\subsection{質問 \& 答え}

\paragraph{質問。} 最後のスライドと最後の言葉は、何をすべきか。

\paragraph{答え。} 最後のスライドは contributions を示し、最後の言葉は、弱々しい礼儀ではなく、本当の締めくくりの身振りを遂行すべきである。

最後のスライドは、質疑のあいだも、人々が部屋を去るあいだも、見えたままである。したがってそれは、仕事の最も強い圧縮形を視界に保っていなければならない。Winston が好むラベルが ``Contributions'' であるのは、それが、何をしなければならなかったか、そして実際に何がなされたかという先のリストを、正確に映し返すからである。

最後の言葉は、また別の技芸である。ジョークは最後にはうまく働きうる、と Winston は言う。そこまで来れば聴衆は話し手に順応しており、そのタイミングを受け取れるからだ。しかし彼は、主たる結びとして ``thank you'' で終えることを勧めない。そうした終わり方は、聴衆がそこにいたのは話を聞く価値があったからではなく、礼儀からだったのだ、という印象を与える。そこで彼は、別の例を探す。政治的な祝祷、儀礼的な解散の定型句、演奏会の慣習、そして最後に、ここで最も勧めたい身振り、すなわち、聴衆とその場に対する敬礼である。

このことも、講義そのものによって実演される。Winston は、ここにいるというそのこと自体によって、プレゼンテーションが重要だという理解を、聴衆はすでに示しているのだと告げ、そのことに対して彼らへ敬礼する。本当の締めの身振りがあって、はじめて最後の丁重な謝意が来る。この順序が大切なのである。

\section{まとめ}

この講義は、意図的な絞り込みと拡大によって進む。まず、
\[
Q=f(K,P,T),
\]
という公式へと絞り込み、ついでただちに \(T\) を縮小することで、その公式の対称性を破る。そこからヒューリスティックへ、つぎに部屋へ、つぎに媒体の選択へ、さらには口頭試問やジョブトークといった特殊な場合へと広がり、そして最後に、Winston's Star、contributions スライド、締めの身振りという、簡潔な要約装置へと、もう一度収束する。

この講義全体を統一しているのは、ただ一つの実践的な問いである。すなわち、時間、注意散漫、身体性、記憶というものとの接触を経ても、アイデアが生き残る最大の可能性を、私たちはどう与えるのか。Winston の答えは、一つの小技ではない。レパートリーである。私たちは動機づけ、約束し、巡回し、柵を立て、句読点を打ち、問いを発し、部屋を選び、媒体を選び、できるときには小道具で議論を運び、そして最後に、何が成し遂げられたのかについて最も強い文を聴衆に残して終えるのである。